\section{Conclusion}
\label{sec:conclusion}

This report presents StepAudio 2.5 as a unified audio-language foundation with three downstream specializations. The shared backbone is learned through a staged multimodal curriculum that aligns speech and text, extends the token interface to audio, scales unified multimodal training, and refines the model with high-quality long-context data. On top of that backbone, the ASR branch turns speech recognition into a particularly favorable application of verifiable multi-token decoding, the TTS branch turns speech generation into a problem of semantic-to-audio alignment strengthened by context-rich supervision and reinforcement learning, and the Realtime branch extends the same foundation to low-latency spoken dialogue with persona stability and paralinguistic sensitivity. Taken together, these capabilities indicate that StepAudio 2.5 is best understood not as a collection of isolated speech endpoints, but as a shared foundation whose recognition, synthesis, and realtime interaction abilities emerge through different optimization and deployment regimes.
